\documentclass[10pt,twocolumn,letterpaper]{article}
\usepackage[margin=0.75in]{geometry}
\usepackage{times}
\usepackage{booktabs}
\usepackage{graphicx}
\usepackage{amsmath}
\usepackage[hidelinks]{hyperref}
\usepackage{titlesec}
\usepackage{caption}
\captionsetup{font=small,labelfont=bf}
\titleformat{\section}{\normalsize\bfseries}{\thesection.}{0.5em}{}
\titleformat{\subsection}{\small\bfseries}{\thesubsection}{0.5em}{}
\titlespacing*{\section}{0pt}{8pt}{3pt}
\titlespacing*{\subsection}{0pt}{6pt}{2pt}
\setlength{\parindent}{1em}
\setlength{\columnsep}{0.28in}

\begin{document}

\twocolumn[
\begin{center}
{\LARGE\bfseries Does 4-Bit Quantization Change Multi-Turn Jailbreak Robustness?}\par\vspace{0.45em}
{\large A Two-Metric Evaluation of NF4, AWQ, and GPTQ on Llama-3.2-3B and Llama-3.1-8B}\par\vspace{0.9em}
{Parv Bansal}\par\vspace{0.15em}
{AIMS-DTU Research Internship 2026}\par\vspace{0.1em}
{\texttt{parvbansal0011@gmail.com}}\par\vspace{1.1em}
\end{center}
\begin{center}
\begin{minipage}{0.86\textwidth}
\small
\textbf{Abstract.}
Quantization lets large language models run on edge hardware, where real use is
multi-turn. Prior work measured quantization's effect on refusal at a single turn
only and reported degradation. This paper gives the first measurement across a
multi-turn jailbreak, testing the claim on Llama-3.2-3B (FP16, NF4) and
Llama-3.1-8B (FP16, AWQ, GPTQ) with two metrics: a refusal classifier and a
Llama Guard 3 harm judge. Quantization does not degrade safety. On standalone
AdvBench prompts, single-turn refusal is flat or higher under quantization for
both models and all three methods, with no significant comparison. Under CoSafe
multi-turn attacks, precision changes neither refusal (3B 23\% versus 23\%,
McNemar $p=0.62$; 8B 27 to 32\%, Cochran $Q$ $p=0.12$) nor harmful-response rate.
The one significant single-turn effect, 8B refusing more after quantization,
vanishes under the attack. A safety-prompt positive control confirms the harm
metric can detect a real intervention, so the null is a measured absence of effect
rather than a lack of power. The evaluation also shows that refusal-based attack
success overstates harm by an order of magnitude: the CoSafe attack drives 8B to
73\% non-refusal, but only 6\% of those responses are actually harmful.
Single-turn safety evaluation does not predict how quantized models rank under
conversational attack.
\end{minipage}
\end{center}
\vspace{1.3em}
]

\section{Introduction}
Post-training quantization to 4 bits is now the standard way to fit an
instruction-tuned model onto a phone or a laptop. The same deployment target that
forces quantization also forces multi-turn use, because a chat assistant holds a
conversation. Safety evaluation has not followed. Recent work that measures how
quantization changes refusal uses single-turn prompts and reports a drop in
refusal after quantization (Kadadekar, 2026). Whether that drop holds, grows, or
disappears once an attacker has several turns is open, and the authors name
multi-turn as future work.

This paper answers one question: does 4-bit quantization change a model's
multi-turn jailbreak robustness, and how does that compare to its single-turn
effect on the same models and prompts? The study covers two models and three
quantization methods, and reads two behavioral signals, whether the model refused
and whether the response is actually harmful, along with the projection onto the
refusal direction of Arditi et al. (2024).

The reported behavioral degradation does not reproduce, at either horizon. On
standalone AdvBench prompts, quantization leaves single-turn refusal flat or
higher for both models and all three methods. Under a CoSafe coreference attack,
precision does not change refusal either, and all precisions converge to a common
non-refusal rate (Figure~\ref{fig:hero}a). A Llama Guard 3 rescore confirms the
null on a second metric: quantization never raises the rate of actually-harmful
responses (Figure~\ref{fig:hero}b). The only significant behavioral precision
effect anywhere is an increase, where 8B quantization raises refusal on the CoSafe
cold turn, and even that vanishes under the attack.

The practical consequence is a warning about evaluation. A single-turn safety
score does not preserve how FP16 and quantized checkpoints rank under
conversational attack, and on these models it does not surface a behavioral safety
cost from quantization at all. The contributions are:

\begin{itemize}
\setlength\itemsep{1pt}
\item A well-powered non-replication across two Llama models, three 4-bit methods,
two prompt sets (AdvBench and CoSafe), and two metrics (substring refusal and a
Llama Guard harm judge).
\item The first measurement of 4-bit quantization under a multi-turn attack.
Precision changes neither multi-turn refusal nor multi-turn harm, and the one
significant single-turn effect does not survive the attack.
\item Two evaluation cautions. Refusal-based attack success overstates actual harm
by an order of magnitude on CoSafe, and a safety-prompt positive control shows the
harm metric can detect a real intervention.
\end{itemize}

\begin{figure*}[t]
\centering
\includegraphics[width=0.92\textwidth]{fig_hero.pdf}
\caption{The two-metric null in one view. (a) Multi-turn refusal does not vary
with precision (Wilson 95\% intervals overlap). (b) Multi-turn harmful-response
rate is low and does not rise under quantization, yet a safety system prompt drops
it to 0\%, showing the metric can detect a real intervention.}
\label{fig:hero}
\end{figure*}

\section{Background}
\subsection{Quantization and safety}
Kadadekar (2026) reports that INT4 quantization degrades single-turn refusal while
the refusal direction stays geometrically intact, and notes that standard probes
do not explain the behavioral gap. Chhabra and Khalili (2025) find that the
refusal direction survives compression with cosine similarity above 0.99. Both
study single-turn behavior. Neither measures a conversation. This paper extends
the first to the multi-turn setting it names as open.

\subsection{Refusal as a direction}
Arditi et al. (2024) show that refusal in instruction-tuned models is mediated by
a single linear direction in the residual stream. Adding it induces refusal, and
ablating it removes refusal. This study uses their difference-in-means estimator
and confirms the causal layer by ablation rather than by separability.

\subsection{Multi-turn attacks}
STAR (2026) shows that in full-precision models the projection onto the refusal
direction decays across conversation turns, so multi-turn state erodes refusal.
CoSafe (Yu et al., 2024) builds multi-turn attacks whose final turn refers back to
earlier turns by coreference, so the harmful request is only legible in context.
The measurement here is independent of STAR and is used to compare precisions.

\section{Experimental Setup}
\subsection{Models and precisions}
The 3B model is Llama-3.2-3B-Instruct at FP16 and NF4 (bitsandbytes
double-quantized). The 8B model is Llama-3.1-8B-Instruct at FP16, AWQ-INT4, and
GPTQ-INT4, using the public \texttt{hugging-quants} checkpoints. AWQ and GPTQ
layers are not fused, so forward hooks see the residual stream.

\subsection{Refusal direction and causal layer}
The refusal direction is extracted by difference-in-means: the mean activation at
the post-instruction token for harmful prompts minus the same for harmless
prompts, per layer, normalized. Held-out separation accuracy is 100\% for the 3B
direction. The functional layer is chosen by ablation, not by separability. Each
candidate layer's direction is removed from the residual stream, and the layer
whose removal collapses refusal is kept. For 3B this is layer 20 and for 8B layer
12, where ablation drops refusal from 94\% to 0\% while a random-direction control
stays at baseline. The separability-argmax layer differs in both models, 28 for 3B
and 31 for 8B, and does not mediate refusal, so only ablation-confirmed layers are
used.

\subsection{Attacks and prompts}
Multi-turn attacks come from CoSafe, using its safe categories (non-violent
unethical behavior, financial crime, privacy violation, and misinformation), 100
scenarios in total. For each scenario the final unsafe turn is measured two ways.
\textbf{Cold} presents the final turn alone. \textbf{In-context} presents the full
conversation, which is the attack. The two conditions share the same final
request, so their difference isolates the effect of the conversational context.
Standalone single-turn prompts come from AdvBench, 100 instructions presented as
one turn, to test the quantization effect on overtly harmful prompts in the same
setting as prior single-turn work.

\subsection{Metrics}
Refusal is detected by a fixed phrase classifier on greedy generation. To measure
whether a response is actually harmful, every completion is rescored with Llama
Guard 3 (Inan et al., 2023), reporting the harmful-response rate. The projection
of the post-instruction activation onto the refusal direction is recorded as a
mechanism readout. Refusal and harm are paired binary across precision, so
comparisons use McNemar's test per pair and Cochran's $Q$ across the three 8B
precisions. Projection is paired continuous, so it uses the paired $t$-test. Every
run is 100 items, so $n=100$ per condition.

\section{Results and Analysis}
\subsection{Multi-turn refusal does not depend on precision}
Table~\ref{tab:means} gives the rates. For 3B the two precisions are identical at
23\% in-context refusal (McNemar $\chi^2=0.25$, $p=0.62$), and the paired
projection difference is null ($t=0.57$, $p=0.57$). For 8B the three precisions
span 27 to 32\%, and Cochran's $Q$ over the paired refusals is not significant
($Q=4.2$, $p=0.12$; Figure~\ref{fig:hero}a). The non-refusal rate sits near 70\%
for every precision, though most of that is not harmful output.

\begin{table}[t]
\centering\small
\begin{tabular}{llcccc}
\toprule
Model & Prec. & \multicolumn{2}{c}{Refusal} & \multicolumn{2}{c}{Proj.} \\
& & cold & in-ctx & cold & in-ctx \\
\midrule
3B & FP16 & 25\% & 23\% & $+1.46$ & $-0.86$ \\
3B & NF4  & 23\% & 23\% & $+1.33$ & $-0.84$ \\
\midrule
8B & FP16 & 23\% & 27\% & $+0.88$ & $+0.82$ \\
8B & AWQ  & 31\% & 32\% & $+0.80$ & $+0.76$ \\
8B & GPTQ & 34\% & 31\% & $+0.90$ & $+0.89$ \\
\bottomrule
\end{tabular}
\caption{Refusal and mean projection onto the refusal direction, cold (single
turn) and in-context (multi-turn), $n=100$ per cell.}
\label{tab:means}
\end{table}

\subsection{The single-turn gap does not survive the attack}
The cold column tells a different story for 8B. Quantization raises single-turn
refusal, with FP16 at 23\%, AWQ at 31\%, and GPTQ at 34\%. Both pairwise contrasts
against FP16 are significant (McNemar $p=0.043$ and $p=0.010$; Table~\ref{tab:sig}),
and Cochran's $Q$ is significant ($Q=9.7$, $p=0.008$). The same three precisions
are not separable once the conversation is present. The 8 to 11 point spread at a
single turn narrows to 5 points in context and loses significance
(Figure~\ref{fig:coldvsic}). The 3B model shows no precision effect at either
horizon.

\begin{figure}[t]
\centering
\includegraphics[width=0.92\columnwidth]{fig2_cold_vs_ic.pdf}
\caption{Llama-3.1-8B refusal, single turn (cold) versus multi-turn (in-context).
The single-turn precision gap narrows and loses significance in context.}
\label{fig:coldvsic}
\end{figure}

\begin{table}[t]
\centering\small
\begin{tabular}{lll}
\toprule
Comparison & Cold & In-context \\
\midrule
3B FP16 vs NF4 & $p=0.68$ & $p=0.62$ \\
8B FP16 vs AWQ & $p=0.043$ & $p=0.13$ \\
8B FP16 vs GPTQ & $p=0.010$ & $p=0.22$ \\
8B Cochran $Q$ & $p=0.008$ & $p=0.12$ \\
\bottomrule
\end{tabular}
\caption{Refusal significance. The single-turn precision effect on 8B is
significant, and the multi-turn effect is not.}
\label{tab:sig}
\end{table}

\subsection{No degradation on standalone prompts}
On AdvBench single-turn prompts (Table~\ref{tab:advbench}) quantization does not
lower refusal either. The 3B model moves from 76\% to 82\% under NF4 (McNemar
$p=0.15$). The 8B model moves from 95\% to 98\% under AWQ ($p=0.25$) and stays at
95\% under GPTQ ($p=0.62$). No comparison is significant, and every point estimate
is flat or higher under quantization. This is a well-powered non-replication of
the single-turn refusal degradation reported for INT4 quantization. Refusal here
is near ceiling for 8B, so the harm judge addresses the question from the other
side.

\begin{table}[t]
\centering\small
\begin{tabular}{llccc}
\toprule
Model & \multicolumn{3}{c}{Refusal} & McNemar \\
& FP16 & \multicolumn{2}{c}{quantized} & vs FP16 \\
\midrule
3B & 76\% & NF4 82\% & & $p=0.15$ \\
8B & 95\% & AWQ 98\% & GPTQ 95\% & $0.25 / 0.62$ \\
\bottomrule
\end{tabular}
\caption{Single-turn refusal on 100 standalone AdvBench prompts. No precision
comparison is significant.}
\label{tab:advbench}
\end{table}

\subsection{The null holds under a harm judge}
Rescoring every completion with Llama Guard 3, which labels whether the response
itself is harmful, gives the rates in Table~\ref{tab:harm}. Quantization does not
raise the harmful-response rate in any condition. Every quantized run is at or
below its FP16 baseline, and no comparison is significant (McNemar $p \geq 0.48$).
The null survives a second, stricter metric (Figure~\ref{fig:twometric}).

\begin{table}[t]
\centering\small
\begin{tabular}{llcc}
\toprule
Setting & Model & FP16 & quantized \\
\midrule
single & 3B & 1\% & NF4 1\% \\
single & 8B & 4\% & AWQ 2\% / GPTQ 3\% \\
multi  & 3B & 4\% & NF4 3\% \\
multi  & 8B & 6\% & AWQ 4\% / GPTQ 5\% \\
\bottomrule
\end{tabular}
\caption{Harmful-response rate (Llama Guard 3), lower is safer. Quantization never
raises it; all comparisons non-significant.}
\label{tab:harm}
\end{table}

\begin{figure*}[t]
\centering
\includegraphics[width=0.92\textwidth]{fig_twometric.pdf}
\caption{Single-turn versus multi-turn on both metrics for Llama-3.1-8B.
Quantization moves neither refusal (a) nor harmful-response rate (b) at either
horizon.}
\label{fig:twometric}
\end{figure*}

\subsection{The harm metric has power: a positive control}
A null on a rare-event metric invites the objection that nothing could show up. To
rule that out, a safety system prompt is prepended to every FP16 condition, an
intervention that should change behavior. It does, and the harm metric detects it.
Under the multi-turn attack, the system prompt cuts 8B harmful responses from 6\%
to 0\% (McNemar $p=0.041$) and raises multi-turn refusal from 27\% to 38\%
($p=0.006$; Figure~\ref{fig:hero}b). The harm metric is not saturated. It registers
a real safety intervention on the same prompts where quantization registers
nothing. The quantization null is a measured absence of effect, not an absence of
power.

\subsection{Refusal-based attack success overstates harm}
The two metrics diverge sharply. The CoSafe attack takes 8B to a 73\% non-refusal
rate by the substring metric, but only 6\% of those responses are judged harmful
by Llama Guard. The attack mostly produces non-refusing but non-harmful text. On
these prompts, refusal-based attack success overstates actual harm by an order of
magnitude, which is a caution for how multi-turn safety is reported.

\subsection{The direction stays coupled}
If quantization broke the refusal representation, the projection would stop
predicting refusal. It does not. The AUC of projection predicting refusal ranges
from 0.79 to 0.93 across all precisions and conditions
(Figure~\ref{fig:coupling}), and point-biserial correlations run from 0.46 to 0.68
with $p<10^{-6}$. Multi-turn convergence is behavioral saturation across a shared
threshold, not a loss of the mechanism.

\begin{figure}[t]
\centering
\includegraphics[width=0.92\columnwidth]{fig3_coupling.pdf}
\caption{AUC of projection onto the refusal direction predicting refusal. The
direction stays coupled to behavior under every precision.}
\label{fig:coupling}
\end{figure}

\section{Discussion}
\subsection{Why the multi-turn gap closes}
CoSafe attacks hide the harmful request behind coreference, so the model answers
the earlier benign-looking turns and then completes the final one in context. The
attack is strong, taking FP16 from 23\% to 27\% refusal for 8B and driving 3B to
23\%. Once the attack is this effective, a few points of extra caution from
quantization is not enough to change the outcome. A safety margin that exists in
isolation does not imply a safety margin under pressure.

\subsection{Implication for evaluation}
Single-turn refusal is the standard axis for comparing a quantized checkpoint to
its FP16 source. This result shows that axis does not preserve the multi-turn
ordering. For 8B, single-turn evaluation ranks the quantized checkpoints as safer
than FP16, while under attack they are statistically tied. A single-turn number
does not transfer to the conversational setting where the model is deployed.

\subsection{Relation to prior reports}
Kadadekar (2026) reports single-turn degradation on standard harmful prompts. This
study does not reproduce it. Across both prompt sets and both metrics the
behavioral effect of quantization is null-to-positive, not degradation. The Llama
Guard rescore closes the ceiling concern that would otherwise bound a refusal-only
null, because on a metric with room below, and one the positive control shows can
detect a real intervention, quantization still moves nothing.

\section{Limitations and Future Work}
The study covers two model sizes (3B and 8B), one family (Llama), one multi-turn
attack (CoSafe coreference), and three quantization methods. It does not vary
attack type or study training-time quantization. The refusal detector is a phrase
classifier, though the harm judge and the projection readout, which do not use it,
agree with the behavioral results. Both metrics sit near an extreme, with
substring refusal near ceiling and the Llama Guard harm rate near floor, so each
alone has limited room to expose a small effect. The positive control mitigates
this for the harm metric, but the low harm rate partly reflects the CoSafe safe
categories being mild relative to the severe-hazard taxonomy Llama Guard is tuned
for. The result is best read as: on these prompts, actual harmful output is rare at
every precision, not as a guarantee on severe-harm categories. A stricter judge, a
second multi-turn attack, and a larger model family are the direct next steps.

\section{Conclusion}
Across two Llama models, three 4-bit methods, two prompt distributions, and two
metrics, 4-bit quantization does not reduce safety at either a single turn or
across a multi-turn attack. The only significant behavioral precision effect is an
increase, where 8B refuses more on the CoSafe cold turn, and it does not survive
the attack, where all precisions converge near 30\% refusal and 4 to 6\% harmful
output. A safety system prompt, by contrast, does lower multi-turn harm on the
same prompts, so the null is a true absence of effect rather than a lack of
measurement power. Single-turn safety evaluation does not predict how quantized
models rank under conversational attack.

\section*{References}
\footnotesize
Arditi, A., Obeso, O., Syed, A., Paleka, D., Panickssery, N., Gurnee, W., and
Nanda, N. (2024). Refusal in Language Models Is Mediated by a Single Direction.
\textit{NeurIPS}.

Chhabra, A., and Khalili, M. M. (2025). Refusal in Compressed Models via
Mechanistic Interpretability. \textit{arXiv:2504.04215}.

Dettmers, T., Pagnoni, A., Holtzman, A., and Zettlemoyer, L. (2023). QLoRA:
Efficient Finetuning of Quantized LLMs. \textit{NeurIPS}.

Frantar, E., Ashkboos, S., Hoefler, T., and Alistarh, D. (2023). GPTQ: Accurate
Post-Training Quantization for Generative Pre-trained Transformers. \textit{ICLR}.

Inan, H., Upasani, K., Chi, J., Rungta, R., Iyer, K., Mao, Y., et al. (2023).
Llama Guard: LLM-based Input-Output Safeguard for Human-AI Conversations.
\textit{arXiv:2312.06674}.

Kadadekar. (2026). Quality Is Not a Safety Proxy Under Quantization.
\textit{arXiv:2606.10154}.

Lin, J., Tang, J., Tang, H., Yang, S., Dang, X., and Han, S. (2024). AWQ:
Activation-Aware Weight Quantization for LLM Compression and Acceleration.
\textit{MLSys}.

STAR. (2026). State-Dependent Safety Failures in Multi-Turn Language Model
Interaction. \textit{arXiv:2603.15684}.

Yu, E., et al. (2024). CoSafe: Evaluating Large Language Model Safety in Multi-Turn
Dialogue Coreference. \textit{EMNLP}.

Zou, A., Wang, Z., Kolter, J. Z., and Fredrikson, M. (2023). Universal and
Transferable Adversarial Attacks on Aligned Language Models.
\textit{arXiv:2307.15043}.

\end{document}
